%%%%%%%% DriftTriage: concept-drift triage for in-context tabular learners. ICML template (icml2026.sty until the
%%%%%%%% 2027 style is released). Skeleton of 2026-10-03: no experiment has been run; every number is a TODO.
\documentclass{article}

\usepackage{microtype}
\usepackage{graphicx}
\usepackage{booktabs}
\usepackage{hyperref}
\newcommand{\theHalgorithm}{\arabic{algorithm}}

\usepackage{icml2026}          % blind review
% \usepackage[accepted]{icml2026}

\usepackage{amsmath}
\usepackage{amssymb}
\usepackage{mathtools}
\usepackage{amsthm}
\usepackage[capitalize,noabbrev]{cleveref}
\usepackage{algorithm}
\usepackage{algorithmic}
\usepackage[textsize=tiny]{todonotes}

\theoremstyle{plain}
\newtheorem{theorem}{Theorem}
\newtheorem{proposition}{Proposition}
\newtheorem{lemma}{Lemma}
\theoremstyle{definition}
\newtheorem{definition}{Definition}
\newtheorem{assumption}{Assumption}
\newtheorem{observation}{Observation}
\theoremstyle{remark}
\newtheorem{remark}{Remark}

% Shared notation. Vectors bold (\mathbf), matrices calligraphic (\mathcal), names and operators upright.
\newcommand{\vx}{\mathbf{x}}
\newcommand{\vX}{\mathbf{X}}
\newcommand{\vz}{\mathbf{z}}
\newcommand{\vZ}{\mathbf{Z}}
\newcommand{\vw}{\mathbf{w}}
\newcommand{\valpha}{\boldsymbol{\alpha}}
\newcommand{\cC}{\mathcal{C}}      % context (a set of labelled rows)
\newcommand{\cY}{\mathcal{Y}}      % label set
\newcommand{\cH}{\mathcal{H}}      % hypothesis space
\newcommand{\cD}{\mathcal{D}}
\newcommand{\cS}{\mathcal{S}}      % data stream
\newcommand{\E}{\mathbb{E}}
\newcommand{\R}{\mathbb{R}}
\newcommand{\Prob}{\Pr}
\newcommand{\KL}{\mathrm{KL}}
\newcommand{\Regret}{\mathrm{Regret}}
\newcommand{\Div}{\mathrm{D}}
\DeclareMathOperator*{\argmin}{arg\,min}
\DeclareMathOperator*{\argmax}{arg\,max}
\DeclareMathOperator{\AUROC}{AUROC}

\crefname{assumption}{Assumption}{Assumptions}
\crefname{observation}{Observation}{Observations}
\newcommand{\method}{\textsc{DriftTriage}}   % working name

\begin{document}

\twocolumn[
  \icmltitle{Benign, Harmful or Novel? Triaging Concept Drift for In-Context Tabular Learners}
  \icmlsetsymbol{equal}{*}
  \begin{icmlauthorlist}
    \icmlauthor{Anonymous Author(s)}{anon}
  \end{icmlauthorlist}
  \icmlaffiliation{anon}{Anonymous Institution}
  \icmlcorrespondingauthor{Anonymous Author(s)}{anon@anon.org}
  \icmlkeywords{concept drift, tabular foundation models, in-context learning, data streams}
  \vskip 0.3in
]
\printAffiliationsAndNotice{}

\begin{abstract}
Tabular foundation models (TFMs) adapt to a data stream by changing their context rather than their weights. We ask
which concept drifts actually harm an in-context learner and which of them it can see. On controlled streams with
TabPFN, a shift of the input distribution inside the support of the context costs no accuracy even when a two-sample
test separates old from new rows with AUROC 0.95, a shift outside the support costs as much as a change of the
labelling function, and a change of the labelling function leaves the model's confidence untouched. We prove that no
label-free statistic can detect such a real drift, that covered drift is benign, and we propose \method{}, which
triages the error stream of a standard drift detector by coverage and class novelty. In a pre-registered held-out
evaluation it matches the accuracy of a detector-triggered reset, cuts needless resets from 47 to 9 and keeps the
detection of new classes outside the support (AUROC 0.79 versus 0.61), while unsupervised drift detectors fire on
benign drift and miss harmful drift.
\end{abstract}

% =====================================================================================================
\section{Introduction}
\label{sec:intro}

\textbf{Momentum.} Prior-data fitted networks approximate Bayesian posterior prediction in one forward pass
\citep{muller2022pfn,nagler2023statistical}, and TFMs such as TabPFN \citep{hollmann2025tabpfn} and TabICL
\citep{qu2025tabicl} now classify tables without gradient updates. Because they learn from the context, they are a
natural fit for data streams, where the context is a sketch of recent data \citep{lourenco2026streams}, and for
temporal shift \citep{helli2024drift}.

\textbf{Problem.} Streams are non-stationary: their joint distribution changes over time
\citep{gama2014survey,lu2019learning}. Concept drift theory separates \emph{virtual} drift of $P(\vx)$ from
\emph{real} drift of $P(y\mid\vx)$ and treats the appearance of an unseen class as a sudden drift. For an in-context
learner these types have different consequences, but existing TFM work treats drift as one phenomenon and measures
only accuracy after adaptation.

\begin{observation}\label{obs:types}
The type of a drift, not the size of the input shift, decides whether a TFM fails. When $P(\vx)$ is tilted inside the
support of the context, a domain classifier separates old from new rows with AUROC up to 0.95, yet the stale context
loses at most 0.006 accuracy, within seed noise. When the labelling function changes on the same inputs, the loss
equals the fraction of rows whose label changed, within 0.021 (hyperplane rotated by $90^\circ$: 0.51 versus 0.98).
When $P(\vx)$ moves outside the support, the stale context fails as well (periodic boundary translated by twice its
range: 0.52 versus 0.99). (Held-out seeds; the pre-registered bound of 0.005 for the tilt was missed by 0.0006.)
\todo{Figure.}
\end{observation}

\begin{observation}\label{obs:confidence}
Label-free signals see the wrong drifts. Under real drift the TFM's confidence does not move (it changes by at most
0.004; 0.974 at $90^\circ$, where accuracy is 0.51) and the input distribution is unchanged, so neither confidence nor a $P(\vx)$
detector can see it; a $P(\vx)$ detector instead fires on the benign in-support tilt. Uncovered inputs, in contrast,
lower the confidence (0.66) and are far from the context: a nearest-neighbour coverage ratio is $\le$1.07 for every
in-support tilt and $\ge$3.6 for every translation.
\end{observation}

Together these observations split concept drift for in-context learners along two axes: whether the drift harms
the in-context predictor, and whether it is visible without labels. This raises the question: \emph{how should an
in-context learner use the signals it has, label-free coverage and delayed labels, to keep its context when drift is
benign and change it only when drift is harmful?}

\textbf{Approach.} We formalise concept drift for in-context learners (\cref{sec:def}) and decompose the
discrepancy between the context and an incoming window into a density ratio $r(\vx)$ and a likelihood ratio
$v(\vx,y)$, following the sampling-transfer view of I-Div \citep{zhao2024idiv}. In practice the two
signals an in-context learner has are label-free coverage, which reveals uncovered drift, and delayed labels, which
alone can reveal real drift (\cref{prop:ident}). \method{} therefore triages the error stream: errors on uncovered rows
extend the context, errors on rows of a new class are rejected, and only the remaining errors can trigger a reset.

\textbf{Contributions.}
(1) A definition of concept drift for in-context learners that separates benign, harmful and novel drift
(\cref{def:icdrift}).
(2) Observations showing that the drift type, not its magnitude, decides whether a TFM fails, and that TFM
confidence cannot tell the types apart.
(3) \method{}, a training-free triage of the error stream, non-inferior in accuracy to the strongest detector while
avoiding needless resets and preserving novelty detection, with propositions on what is and is not identifiable.
(4) A benchmark of tabular streams with typed drift events and pre-registered held-out evaluation.

% =====================================================================================================
\section{Concept Drift for In-Context Learners}
\label{sec:def}

Let a stream $\cS_{0,t}=\{(\vx_i,y_i)\}$ arrive in windows $W_1,W_2,\dots$ and let $\cC_t$ be the context of a frozen
TFM $f_{\Theta}(\cdot\,;\cC_t)$ at time $t$, a subset of the labelled rows seen so far. Following
\citet{lu2019learning}, concept drift occurs at $t$ if $P_t(\vx,y)\neq P_{t+1}(\vx,y)$, with
$P_t(\vx,y)=P_t(\vx)\,P_t(y\mid\vx)$.

\begin{definition}[In-context concept drift]\label{def:icdrift}
Let $P_{\cC}$ be the distribution of the context and $P_{W}$ that of an incoming window, with label sets $\cY_{\cC}$
and $\cY_W$. Write $r(\vx)=P_W(\vx)/P_{\cC}(\vx)$ and $v(\vx,y)=P_W(y\mid\vx)/P_{\cC}(y\mid\vx)$ for
$y\in\cY_{\cC}$. The drift from $\cC$ to $W$ is
\begin{itemize}
  \item \emph{covered} (benign virtual) if $v\equiv1$ and $r$ is bounded on the support of $P_W$;
  \item \emph{uncovered} (coverage drift) if $v\equiv1$ but $P_W$ puts mass where $P_{\cC}$ has none ($r$ unbounded);
  \item \emph{real} if $v\not\equiv1$ on a set of positive $P_W$-measure;
  \item \emph{novel} if $P_W(y\notin\cY_{\cC})>0$.
\end{itemize}
\end{definition}

Covered drift is harmless to an in-context learner and uncovered drift is visible without labels; real drift is
harmful and invisible without labels; novel rows must be rejected rather than used to change the context
(Observations~\ref{obs:types},~\ref{obs:confidence}).

\begin{remark}[Identifiability]
Real drift with an unchanged $P(\vx)$ cannot be detected from unlabelled rows alone. We therefore work in the standard
prequential setting with delayed labels \citep{gama2014survey}: the labels of window $W_t$ become available after its
rows are predicted, and $v$ is estimated on the most recent labelled window, while $r$ and novelty are estimated on
the current, unlabelled one. \todo{Check whether I-Div's likelihood-ratio step gives a label-free proxy under an
extra assumption (e.g.\ label shift).}
\end{remark}

% =====================================================================================================
\section{Method}
\label{sec:method}

\method{} runs on a frozen TFM with a context of at most $M$ rows, batches of $B$ rows and labels one batch late.

\textbf{Coverage.} For each incoming (unlabelled) row, the coverage ratio is its nearest-neighbour distance to the
context divided by the mean leave-one-out nearest-neighbour distance within the context; a row is \emph{uncovered} if
the ratio exceeds $\tau=3$. Uncovered rows, once labelled, enter the context with priority: when the context is full,
covered rows are evicted first.

\textbf{Filtered detection.} The DDM rule \citep{gama2014survey} runs only on the errors of rows that were covered and
whose class is already in the context. Errors on uncovered rows signal coverage drift and errors on rows of a new
class signal novelty; neither resets the context. A detection resets the context to the newly labelled rows.

\textbf{Rejection.} Every incoming row gets a rejection score, the larger of one minus the TFM confidence and the
excess coverage ratio $\max(0,\text{ratio}/\tau-1)$.

The only constant chosen on development streams is $\tau$; the DDM rule keeps its standard thresholds. The design
history (two rejected variants based on context-swap tests, and two rejected refinements) is reported in the appendix.
(\cref{app:history}).

\section{Theory}
\label{sec:theory}

Let $\cC$ be the context, $W$ an incoming window, $h=f_{\Theta}(\cdot\,;\cC)$ the in-context predictor and
$R_P(h)=\Prob_{(\vx,y)\sim P}(h(\vx)\neq y)$ its error under $P$.

\begin{proposition}[Real drift is invisible without labels]\label{prop:ident}
Let $T$ be any statistic of the context and the unlabelled rows of $W$. If two window distributions $P_W$ and $P'_W$
have the same input marginal, $P_W(\vx)=P'_W(\vx)$, then $T$ has the same law under both. In particular, no label-free
detector, including the confidence of $h$ and any test of $P(\vx)$, has power above its level against a real drift that
leaves $P(\vx)$ unchanged.
\end{proposition}

\begin{proof}
$T$ is a measurable function of $(\cC,\vX_W)$, and the law of $\vX_W$ is the product of the input marginal, which is
the same under $P_W$ and $P'_W$; $\cC$ is fixed. Hence the laws of $T$ coincide.
\end{proof}

\begin{proposition}[Covered drift is benign]\label{prop:covered}
If the drift from $P_{\cC}$ to $P_W$ is covered in the sense of \cref{def:icdrift}, with $v\equiv1$ and
$r\le\bar r$ on the support of $P_W$, then $R_{P_W}(h)\le\bar r\,R_{P_{\cC}}(h)$.
\end{proposition}

\begin{proof}
With $v\equiv1$ the conditional label law is the same, so the error of $h$ at $\vx$ is the same function
$e(\vx)=\Prob(h(\vx)\neq y\mid\vx)$ under both, and $R_{P_W}(h)=\E_{P_{\cC}}[r(\vx)e(\vx)]\le\bar r\,R_{P_{\cC}}(h)$.
\end{proof}

\Cref{prop:covered} explains the first half of \cref{obs:types}: a TFM that is accurate on its context stays accurate
under any tilt with a bounded density ratio, however large the shift measured by a two-sample test. When $P_W$ puts mass
where $P_{\cC}$ has none, $\bar r=\infty$ and nothing is guaranteed; under a kernel posterior predictive the confidence
then reverts to the prior far from the context, which makes uncovered drift visible without labels.
\Cref{prop:ident} explains the second half of \cref{obs:confidence}, and it is why \method{} detects real drift from
delayed labels only.

\begin{proposition}[Filtered detection]\label{prop:filter}
Suppose that between two resets there is no real drift on covered rows of known classes. Then the errors of $h$ on
the rows that \method{} passes to the detector (covered rows of known classes) are independent with mean at most
$\bar r\,R_{P_{\cC}}(h)$, whatever coverage drift or novel classes occur meanwhile, so the false-alarm guarantee of
the detector \citep{gama2014survey} applies to them.\todo{State the DDM false-alarm property precisely.}
\end{proposition}

% =====================================================================================================
\section{Experiments}
\label{sec:exp}

\textbf{Setup.} Streams of two families with typed drift events every 3000 rows, two cycles over \{tilt, translate,
real, novel\} in random order: \emph{sine2d} (periodic boundary; the novel class is a band outside the support) and
\emph{lin} (four linear classes; the novel class takes over part of the support). Batches of 100 rows, labels one
batch late, a context of 1000 rows, TabPFN v2. Baselines: a FIFO context, the dual memory of
\citet{lourenco2026streams}, a context reset by DDM \citep{gama2014survey} on the error stream, and a reset by an
unsupervised $P(\vx)$ detector (two-sample classifier). All design choices were made on development seeds 0--4; the
hypotheses below were written down before the single run on seeds 10--19.

\begin{table}[t]
  \centering
  \caption{Held-out streams (seeds 10--19). Acc: prequential accuracy; Real: accuracy in the 10 batches after a real
  drift; Needless: resets in the 10 batches after a tilt, a translation or a novel class (sum over 20 streams); Nov:
  AUROC for rows of a new class before their labels arrive (sine2d).}
  \label{tab:stream}
  \small
  \setlength{\tabcolsep}{3pt}
  \begin{tabular}{lccccc}
    \toprule
    Policy & Acc sine2d & Acc lin & Real & Needless & Nov \\
    \midrule
    FIFO            & 0.968 & 0.952 & 0.731 & 0  & 0.762 \\
    Dual memory     & 0.969 & 0.948 & 0.738 & 0  & 0.763 \\
    $P(\vx)$ reset  & 0.967 & 0.950 & 0.731 & 69 & 0.729 \\
    DDM reset       & 0.978 & 0.969 & 0.913 & 47 & 0.606 \\
    \method{}       & 0.980 & 0.969 & 0.912 & 9  & 0.786 \\
    \bottomrule
  \end{tabular}
\end{table}

\textbf{Results} (\cref{tab:stream}). Detection-driven resets dominate: after a real drift, FIFO and the dual memory
recover to 0.73--0.74 within ten batches, a DDM reset to 0.91, because a TFM relearns a concept from a few hundred
rows. The unsupervised $P(\vx)$ detector resets 68 times after benign tilts and translations and never after a real
drift, as \cref{prop:ident} predicts. \method{} is non-inferior to the DDM reset in accuracy (sine2d $+0.17$ points,
better on 10 of 10 streams; lin $-0.06$ points), never resets after a coverage drift, makes 9 instead of 47 needless
resets, and keeps the confidence-based detection of a new class outside the support (AUROC 0.786 versus 0.606,
Wilcoxon $p=0.001$). It does not improve on DDM in accuracy, and when a new class appears inside the support (lin), the
new class is also a real drift for the region it takes over; both reset policies then lose novelty detection (0.66 and
0.64 versus 0.87 for FIFO). All four pre-registered stream hypotheses held; of the observation checks, the bound on the
accuracy loss under in-support tilt was missed by 0.0006 (\cref{obs:types}).

\section{Related Work}
Concept drift \citep{gama2014survey,lu2019learning}; concept drift for new model families, e.g.\ multi-modal LLMs
\citep{yang2025adapting,yang2025walking,yang2026turning}; TFMs under shift \citep{helli2024drift,lourenco2026streams}.
\textbf{Difference to the competitors.} Drift-Resilient TabPFN retrains the prior to extrapolate temporal shift and
\citet{lourenco2026streams} manage the context with a fixed FIFO memory; both target accuracy after adaptation and
treat drift as one phenomenon. \method{} is training-free and decides the drift type before acting, including the
rejection of unseen classes.


\appendix
\section{Design History and Pre-Registration}
\label{app:history}
All choices were made on development seeds 0--4 and frozen before the single run on held-out seeds 10--19.
\textbf{v0} compared the current context with a context of the two previous batches on the newest labelled batch
(context-swap test); it detected real drift about three batches late, because the comparison context still held
pre-drift rows. \textbf{v1} cross-fitted the newest batch instead; with 50 rows per fold it could not learn complex
boundaries and missed real drift (0 resets on one stream). \textbf{v2} (the evaluated method) filters the error stream
of a standard DDM detector. Two refinements were rejected on development streams: counting a class as known only after
30 context rows (cascades of resets after every reset) and a refractory period with local pruning of rows taken over by
a new class (lower accuracy; the refractory period also made DDM worse). The hypotheses of \cref{sec:exp}, including
the observation checks, were written down before the held-out run; all four stream hypotheses held and one observation
bound was missed by 0.0006.

\bibliography{refs}
\bibliographystyle{icml2026}

\end{document}
